% Options for packages loaded elsewhere
\PassOptionsToPackage{unicode}{hyperref}
\PassOptionsToPackage{hyphens}{url}
\documentclass[
]{article}
\usepackage{xcolor}
\usepackage{amsmath,amssymb}
\setcounter{secnumdepth}{-\maxdimen} % remove section numbering
\usepackage{iftex}
\ifPDFTeX
  \usepackage[T1]{fontenc}
  \usepackage[utf8]{inputenc}
  \usepackage{textcomp} % provide euro and other symbols
\else % if luatex or xetex
  \usepackage{unicode-math} % this also loads fontspec
  \defaultfontfeatures{Scale=MatchLowercase}
  \defaultfontfeatures[\rmfamily]{Ligatures=TeX,Scale=1}
\fi
\usepackage{lmodern}
\ifPDFTeX\else
  % xetex/luatex font selection
\fi
% Use upquote if available, for straight quotes in verbatim environments
\IfFileExists{upquote.sty}{\usepackage{upquote}}{}
\IfFileExists{microtype.sty}{% use microtype if available
  \usepackage[]{microtype}
  \UseMicrotypeSet[protrusion]{basicmath} % disable protrusion for tt fonts
}{}
\makeatletter
\@ifundefined{KOMAClassName}{% if non-KOMA class
  \IfFileExists{parskip.sty}{%
    \usepackage{parskip}
  }{% else
    \setlength{\parindent}{0pt}
    \setlength{\parskip}{6pt plus 2pt minus 1pt}}
}{% if KOMA class
  \KOMAoptions{parskip=half}}
\makeatother
\usepackage{color}
\usepackage{fancyvrb}
\newcommand{\VerbBar}{|}
\newcommand{\VERB}{\Verb[commandchars=\\\{\}]}
\DefineVerbatimEnvironment{Highlighting}{Verbatim}{commandchars=\\\{\}}
% Add ',fontsize=\small' for more characters per line
\newenvironment{Shaded}{}{}
\newcommand{\AlertTok}[1]{\textcolor[rgb]{1.00,0.00,0.00}{\textbf{#1}}}
\newcommand{\AnnotationTok}[1]{\textcolor[rgb]{0.38,0.63,0.69}{\textbf{\textit{#1}}}}
\newcommand{\AttributeTok}[1]{\textcolor[rgb]{0.49,0.56,0.16}{#1}}
\newcommand{\BaseNTok}[1]{\textcolor[rgb]{0.25,0.63,0.44}{#1}}
\newcommand{\BuiltInTok}[1]{\textcolor[rgb]{0.00,0.50,0.00}{#1}}
\newcommand{\CharTok}[1]{\textcolor[rgb]{0.25,0.44,0.63}{#1}}
\newcommand{\CommentTok}[1]{\textcolor[rgb]{0.38,0.63,0.69}{\textit{#1}}}
\newcommand{\CommentVarTok}[1]{\textcolor[rgb]{0.38,0.63,0.69}{\textbf{\textit{#1}}}}
\newcommand{\ConstantTok}[1]{\textcolor[rgb]{0.53,0.00,0.00}{#1}}
\newcommand{\ControlFlowTok}[1]{\textcolor[rgb]{0.00,0.44,0.13}{\textbf{#1}}}
\newcommand{\DataTypeTok}[1]{\textcolor[rgb]{0.56,0.13,0.00}{#1}}
\newcommand{\DecValTok}[1]{\textcolor[rgb]{0.25,0.63,0.44}{#1}}
\newcommand{\DocumentationTok}[1]{\textcolor[rgb]{0.73,0.13,0.13}{\textit{#1}}}
\newcommand{\ErrorTok}[1]{\textcolor[rgb]{1.00,0.00,0.00}{\textbf{#1}}}
\newcommand{\ExtensionTok}[1]{#1}
\newcommand{\FloatTok}[1]{\textcolor[rgb]{0.25,0.63,0.44}{#1}}
\newcommand{\FunctionTok}[1]{\textcolor[rgb]{0.02,0.16,0.49}{#1}}
\newcommand{\ImportTok}[1]{\textcolor[rgb]{0.00,0.50,0.00}{\textbf{#1}}}
\newcommand{\InformationTok}[1]{\textcolor[rgb]{0.38,0.63,0.69}{\textbf{\textit{#1}}}}
\newcommand{\KeywordTok}[1]{\textcolor[rgb]{0.00,0.44,0.13}{\textbf{#1}}}
\newcommand{\NormalTok}[1]{#1}
\newcommand{\OperatorTok}[1]{\textcolor[rgb]{0.40,0.40,0.40}{#1}}
\newcommand{\OtherTok}[1]{\textcolor[rgb]{0.00,0.44,0.13}{#1}}
\newcommand{\PreprocessorTok}[1]{\textcolor[rgb]{0.74,0.48,0.00}{#1}}
\newcommand{\RegionMarkerTok}[1]{#1}
\newcommand{\SpecialCharTok}[1]{\textcolor[rgb]{0.25,0.44,0.63}{#1}}
\newcommand{\SpecialStringTok}[1]{\textcolor[rgb]{0.73,0.40,0.53}{#1}}
\newcommand{\StringTok}[1]{\textcolor[rgb]{0.25,0.44,0.63}{#1}}
\newcommand{\VariableTok}[1]{\textcolor[rgb]{0.10,0.09,0.49}{#1}}
\newcommand{\VerbatimStringTok}[1]{\textcolor[rgb]{0.25,0.44,0.63}{#1}}
\newcommand{\WarningTok}[1]{\textcolor[rgb]{0.38,0.63,0.69}{\textbf{\textit{#1}}}}
\usepackage{longtable,booktabs,array}
\newcounter{none} % for unnumbered tables
\usepackage{calc} % for calculating minipage widths
% Correct order of tables after \paragraph or \subparagraph
\usepackage{etoolbox}
\makeatletter
\patchcmd\longtable{\par}{\if@noskipsec\mbox{}\fi\par}{}{}
\makeatother
% Allow footnotes in longtable head/foot
\IfFileExists{footnotehyper.sty}{\usepackage{footnotehyper}}{\usepackage{footnote}}
\makesavenoteenv{longtable}
\setlength{\emergencystretch}{3em} % prevent overfull lines
\providecommand{\tightlist}{%
  \setlength{\itemsep}{0pt}\setlength{\parskip}{0pt}}
\usepackage{bookmark}
\IfFileExists{xurl.sty}{\usepackage{xurl}}{} % add URL line breaks if available
\urlstyle{same}
\hypersetup{
  hidelinks,
  pdfcreator={LaTeX via pandoc}}

\author{}
\date{}

\begin{document}

\section{Security Audit: Lux Governance
Contracts}\label{security-audit-lux-governance-contracts}

\textbf{Auditor:} Trail of Bits-level Security Review \textbf{Date:}
2026-01-30 \textbf{Scope:} \texttt{/contracts/governance/} - Core
governance system \textbf{Contracts Reviewed:}

\begin{itemize}
\tightlist
\item
  Stake.sol
\item
  Charter.sol
\item
  Council.sol
\item
  Committee.sol
\item
  Governance.sol (re-exports only)
\item
  MultiDAOGovernor.sol
\item
  Secretariat.sol
\item
  veto/Veto.sol
\item
  veto/Sanction.sol
\item
  voting/VotingWeightLRC20.sol
\item
  voting/VoteTrackerLRC20.sol
\end{itemize}

\begin{center}\rule{0.5\linewidth}{0.5pt}\end{center}

\subsection{Executive Summary}\label{executive-summary}

The Lux governance system implements a modular DAO governance
architecture with Council-Charter separation, multi-DAO coordination,
and parent-child veto mechanics. This audit identified \textbf{4
critical}, \textbf{6 high}, \textbf{9 medium}, and \textbf{12 low}
severity issues across the contract suite.

{\def\LTcaptype{none} % do not increment counter
\begin{longtable}[]{@{}ll@{}}
\toprule\noalign{}
Severity & Count \\
\midrule\noalign{}
\endhead
\bottomrule\noalign{}
\endlastfoot
Critical & 4 \\
High & 6 \\
Medium & 9 \\
Low & 12 \\
\end{longtable}
}

\begin{center}\rule{0.5\linewidth}{0.5pt}\end{center}

\subsection{Critical Findings}\label{critical-findings}

\subsubsection{C-01: Flash Loan Voting Attack in
Charter.sol}\label{c-01-flash-loan-voting-attack-in-chartersol}

\textbf{Location:} \texttt{Charter.sol:395-458} (\texttt{castVote})
\textbf{Severity:} Critical \textbf{Impact:} Complete governance
takeover

\textbf{Description:} The voting weight is calculated at
\texttt{proposal.votingStartTimestamp} using \texttt{getPastVotes}, but
this does not prevent flash loan attacks within the same block as
proposal creation. An attacker can:

\begin{enumerate}
\def\labelenumi{\arabic{enumi}.}
\tightlist
\item
  Monitor mempool for \texttt{submitProposal} transactions
\item
  Frontrun with flash loan to acquire tokens
\item
  Delegate tokens to themselves (in same transaction)
\item
  Wait for proposal initialization (same block,
  \texttt{votingStartTimestamp\ =\ block.timestamp})
\item
  Vote with massive weight
\end{enumerate}

\textbf{Vulnerable Code:}

\begin{Shaded}
\begin{Highlighting}[]
\NormalTok{// Charter.sol:433}
\NormalTok{(uint256 weight, bytes memory processedData) = IVotingWeight(config.votingWeight)}
\NormalTok{    .calculateWeight(voter, proposal.votingStartTimestamp, configData.voteData);}
\end{Highlighting}
\end{Shaded}

\textbf{Root Cause:} The \texttt{votingStartTimestamp} is set to
\texttt{block.timestamp} in \texttt{initializeProposal()}, allowing
same-block voting with manipulated balances.

\textbf{Recommendation:} Add a minimum voting delay (e.g., 1 block)
between proposal creation and voting start:

\begin{Shaded}
\begin{Highlighting}[]
\NormalTok{proposal.votingStartTimestamp = uint48(block.timestamp + MIN\_VOTING\_DELAY);}
\end{Highlighting}
\end{Shaded}

\begin{center}\rule{0.5\linewidth}{0.5pt}\end{center}

\subsubsection{C-02: Quorum Manipulation via Checkpoint Timing in
VotingWeightLRC20.sol}\label{c-02-quorum-manipulation-via-checkpoint-timing-in-votingweightlrc20sol}

\textbf{Location:} \texttt{VotingWeightLRC20.sol:104-113}
(\texttt{calculateWeight}) \textbf{Severity:} Critical \textbf{Impact:}
Artificial quorum satisfaction

\textbf{Description:} An attacker can manipulate quorum by acquiring
tokens just before proposal creation, voting, then transferring tokens
to other controlled addresses that also vote. Since weight is calculated
at \texttt{votingStartTimestamp}, the same tokens can be used to inflate
quorum through circular transfers completed before the snapshot.

\textbf{Attack Vector:}

\begin{enumerate}
\def\labelenumi{\arabic{enumi}.}
\tightlist
\item
  Attacker creates N addresses
\item
  Acquires large token position
\item
  Creates checkpoint by delegating at each address sequentially
\item
  Creates proposal (snapshot taken)
\item
  Each address has the full balance at snapshot time if delegations were
  orchestrated correctly
\end{enumerate}

\textbf{Recommendation:} Use block-based snapshots instead of
timestamp-based, and ensure delegation changes require minimum settling
period.

\begin{center}\rule{0.5\linewidth}{0.5pt}\end{center}

\subsubsection{C-03: Cross-DAO Vote Replay in
MultiDAOGovernor.sol}\label{c-03-cross-dao-vote-replay-in-multidaogovernorsol}

\textbf{Location:} \texttt{MultiDAOGovernor.sol:427-461}
(\texttt{castVote}) \textbf{Severity:} Critical \textbf{Impact:} Vote
weight double-counting across DAOs

\textbf{Description:} The \texttt{castVote} function allows voting on a
per-DAO basis with the same token balance. A user with 1000 tokens can
vote on 5 DAOs with 1000 weight each (total 5000 weight influence from
1000 tokens).

\textbf{Vulnerable Code:}

\begin{Shaded}
\begin{Highlighting}[]
\NormalTok{// MultiDAOGovernor.sol:446}
\NormalTok{uint256 weight = governanceToken.getPastVotes(msg.sender, proposal.snapshotBlock);}
\NormalTok{// No check that weight hasn\textquotesingle{}t been used in another DAO for same proposal}
\end{Highlighting}
\end{Shaded}

\textbf{Design Flaw:} The delegation system (\texttt{delegateMultiple}
with weights) suggests intended vote splitting, but \texttt{castVote}
ignores the \texttt{daoWeight} mapping entirely.

\textbf{Recommendation:} Either:

\begin{enumerate}
\def\labelenumi{\arabic{enumi}.}
\tightlist
\item
  Enforce vote weight splitting when voting across multiple DAOs
\item
  Track cumulative weight used per proposal per voter globally
\item
  Require explicit weight allocation at vote time
\end{enumerate}

\begin{center}\rule{0.5\linewidth}{0.5pt}\end{center}

\subsubsection{C-04: Veto Race Condition in
Veto.sol}\label{c-04-veto-race-condition-in-vetosol}

\textbf{Location:} \texttt{Veto.sol:163-192} (\texttt{castVetoVote})
\textbf{Severity:} Critical \textbf{Impact:} Veto bypass through
proposal expiration manipulation

\textbf{Description:} The veto proposal period check and vote recording
are not atomic. An attacker can exploit the window where:

\begin{enumerate}
\def\labelenumi{\arabic{enumi}.}
\tightlist
\item
  Veto proposal is near expiration
\item
  Attacker initiates a new proposal (resetting the counter)
\item
  Legitimate veto votes are lost
\end{enumerate}

\textbf{Vulnerable Code:}

\begin{Shaded}
\begin{Highlighting}[]
\NormalTok{// Veto.sol:167{-}173}
\NormalTok{bool proposalExpired = $.vetoProposalCreated == 0 ||}
\NormalTok{    block.timestamp \textgreater{} $.vetoProposalCreated + $.vetoProposalPeriod;}

\NormalTok{if (proposalExpired) \{}
\NormalTok{    \_initializeVetoProposal();  // Resets vetoProposalVoteCount to 0}
\NormalTok{\}}
\end{Highlighting}
\end{Shaded}

Any user (including attackers) can trigger proposal expiration reset by
calling \texttt{castVetoVote}, wiping accumulated legitimate votes.

\textbf{Recommendation:} Add cooldown period before new veto proposals
can be created, or require admin action to initialize new veto
proposals.

\begin{center}\rule{0.5\linewidth}{0.5pt}\end{center}

\subsection{High Severity Findings}\label{high-severity-findings}

\subsubsection{H-01: Delegation Bypass in Charter Light Account
Resolution}\label{h-01-delegation-bypass-in-charter-light-account-resolution}

\textbf{Location:} \texttt{Charter.sol:508-530}
(\texttt{\_resolveVoter}) \textbf{Severity:} High \textbf{Impact:} Vote
weight theft

\textbf{Description:} The light account resolution logic contains a
logical flaw. When \texttt{lightAccountIndex\ \textgreater{}\ 0}:

\begin{Shaded}
\begin{Highlighting}[]
\NormalTok{if (lightAccount == sender) \{}
\NormalTok{    return sender;  // This condition is backwards}
\NormalTok{\}}
\end{Highlighting}
\end{Shaded}

The check should verify that calling
\texttt{getAddress(sender,\ lightAccountIndex)} returns an account the
sender controls, but instead returns \texttt{sender} unconditionally
when light account equals sender (which would be a collision).

\textbf{Recommendation:} Verify the light account factory correctly maps
the index to an owned account, and return the actual owner from the
light account.

\begin{center}\rule{0.5\linewidth}{0.5pt}\end{center}

\subsubsection{H-02: Missing Proposal Count Bounds in
Council.sol}\label{h-02-missing-proposal-count-bounds-in-councilsol}

\textbf{Location:} \texttt{Council.sol:293-339}
(\texttt{submitProposal}) \textbf{Severity:} High \textbf{Impact:}
Denial of service via proposal spam

\textbf{Description:} No limit on \texttt{totalProposalCount}. An
attacker with proposer rights can spam unlimited proposals, consuming
storage and making proposal enumeration impossible.

\begin{Shaded}
\begin{Highlighting}[]
\NormalTok{// No maximum check}
\NormalTok{$.totalProposalCount++;}
\end{Highlighting}
\end{Shaded}

\textbf{Recommendation:} Add maximum proposal limit or require proposal
deposits that are returned on execution/expiration.

\begin{center}\rule{0.5\linewidth}{0.5pt}\end{center}

\subsubsection{H-03: Reentrancy in Committee Vote
Execution}\label{h-03-reentrancy-in-committee-vote-execution}

\textbf{Location:} \texttt{Committee.sol:161-171} (\texttt{execute})
\textbf{Severity:} High \textbf{Impact:} Potential state manipulation

\textbf{Description:} While \texttt{ReentrancyGuard} is inherited, the
\texttt{execute} function only marks \texttt{proposal.executed\ =\ true}
after the state check. A malicious executor could potentially exploit
state changes between external calls (if any are added).

More critically, the contract emits an event but performs no actual
execution logic - the \texttt{execute} function is a no-op beyond
marking executed, suggesting incomplete implementation.

\textbf{Recommendation:} Complete the execution logic or document that
Committee is signaling-only. Add explicit CEI pattern if execution is
added.

\begin{center}\rule{0.5\linewidth}{0.5pt}\end{center}

\subsubsection{H-04: Soulbound Mode Toggle Attack in
Stake.sol}\label{h-04-soulbound-mode-toggle-attack-in-stakesol}

\textbf{Location:} \texttt{Stake.sol:119-122} (\texttt{setSoulbound})
\textbf{Severity:} High \textbf{Impact:} Governance token liquidity
manipulation

\textbf{Description:} The owner can toggle soulbound mode at will. This
enables a governance attack:

\begin{enumerate}
\def\labelenumi{\arabic{enumi}.}
\tightlist
\item
  DAO votes to make token soulbound (locking voting power)
\item
  Malicious owner disables soulbound before attack
\item
  Attacker acquires tokens on market
\item
  Attacker votes
\item
  Owner re-enables soulbound
\end{enumerate}

\textbf{Recommendation:} Require timelock for soulbound mode changes, or
make soulbound mode immutable after initial setting.

\begin{center}\rule{0.5\linewidth}{0.5pt}\end{center}

\subsubsection{H-05: Bridge Tally Manipulation in
MultiDAOGovernor.sol}\label{h-05-bridge-tally-manipulation-in-multidaogovernorsol}

\textbf{Location:} \texttt{MultiDAOGovernor.sol:466-493}
(\texttt{receiveTally}) \textbf{Severity:} High \textbf{Impact:}
Cross-chain vote manipulation

\textbf{Description:} The \texttt{receiveTally} function blindly accepts
any tally from the \texttt{BRIDGE\_ROLE} without verifying:

\begin{itemize}
\tightlist
\item
  The source chain
\item
  The attestation data (currently ignored)
\item
  Whether votes were already recorded locally
\end{itemize}

\begin{Shaded}
\begin{Highlighting}[]
\NormalTok{// Attestation parameter is completely ignored}
\NormalTok{bytes calldata /* attestation */}
\end{Highlighting}
\end{Shaded}

An attacker controlling the bridge role can inject arbitrary vote
tallies.

\textbf{Recommendation:} Implement attestation verification using a
threshold signature scheme or merkle proof.

\begin{center}\rule{0.5\linewidth}{0.5pt}\end{center}

\subsubsection{H-06: Sanction Guard Can Be Removed by Child
DAO}\label{h-06-sanction-guard-can-be-removed-by-child-dao}

\textbf{Location:} \texttt{Sanction.sol:80-89} (\texttt{initialize}) and
\texttt{Secretariat.sol:93-96} (\texttt{setGuard}) \textbf{Severity:}
High \textbf{Impact:} Veto mechanism bypass

\textbf{Description:} The Sanction guard is set via
\texttt{Secretariat.setGuard()} which requires \texttt{onlyVault}.
However, if the child DAO\textquotesingle s Safe executes a transaction
to remove the guard, the veto mechanism becomes ineffective.

\textbf{Attack Path:}

\begin{enumerate}
\def\labelenumi{\arabic{enumi}.}
\tightlist
\item
  Child DAO proposes removing the Sanction guard
\item
  Proposal passes before parent can veto
\item
  Guard removed, veto mechanism bypassed
\end{enumerate}

\textbf{Recommendation:} Implement immutable guard attachment or require
parent DAO approval for guard changes.

\begin{center}\rule{0.5\linewidth}{0.5pt}\end{center}

\subsection{Medium Severity Findings}\label{medium-severity-findings}

\subsubsection{M-01: Quorum Denominator Missing in
Committee.sol}\label{m-01-quorum-denominator-missing-in-committeesol}

\textbf{Location:} \texttt{Committee.sol:76,\ 167} \textbf{Severity:}
Medium \textbf{Impact:} Quorum check never enforced

\textbf{Description:} The \texttt{quorumPercentage} storage variable is
set but never used in voting logic:

\begin{Shaded}
\begin{Highlighting}[]
\NormalTok{require(proposal.forVotes \textgreater{} proposal.againstVotes, "Proposal defeated");}
\NormalTok{// Missing: quorum check against total votes}
\end{Highlighting}
\end{Shaded}

\textbf{Recommendation:} Implement quorum check in \texttt{execute()}:

\begin{Shaded}
\begin{Highlighting}[]
\NormalTok{require(totalVotes * 100 \textgreater{}= totalSupply * quorumPercentage, "Quorum not reached");}
\end{Highlighting}
\end{Shaded}

\begin{center}\rule{0.5\linewidth}{0.5pt}\end{center}

\subsubsection{M-02: Timestamp Manipulation in Charter Voting
Period}\label{m-02-timestamp-manipulation-in-charter-voting-period}

\textbf{Location:} \texttt{Charter.sol:412-419} \textbf{Severity:}
Medium \textbf{Impact:} Vote timing manipulation

\textbf{Description:} Voting end is checked against
\texttt{block.timestamp} which can be manipulated by miners within
\textasciitilde900 seconds on most chains. This allows miners to:

\begin{itemize}
\tightlist
\item
  Accept votes slightly after the deadline
\item
  Reject votes slightly before the deadline
\end{itemize}

\textbf{Recommendation:} Use block numbers for timing-critical
operations or add tolerance margin.

\begin{center}\rule{0.5\linewidth}{0.5pt}\end{center}

\subsubsection{M-03: DID Immutability in
Stake.sol}\label{m-03-did-immutability-in-stakesol}

\textbf{Location:} \texttt{Stake.sol:140-146} (\texttt{linkDID})
\textbf{Severity:} Medium \textbf{Impact:} Permanent DID lock

\textbf{Description:} Once a DID is linked, it cannot be updated. If a
DID document changes or a user loses control of their DID, they cannot
update it.

\begin{Shaded}
\begin{Highlighting}[]
\NormalTok{if (bytes(did[msg.sender]).length \textgreater{} 0) \{}
\NormalTok{    revert DIDAlreadyLinked();}
\NormalTok{\}}
\end{Highlighting}
\end{Shaded}

\textbf{Recommendation:} Add an \texttt{updateDID()} function with
appropriate timelock or governance controls.

\begin{center}\rule{0.5\linewidth}{0.5pt}\end{center}

\subsubsection{M-04: Veto Proposal Period Not Configurable
Post-Deployment}\label{m-04-veto-proposal-period-not-configurable-post-deployment}

\textbf{Location:} \texttt{Veto.sol:91-104} (\texttt{initialize})
\textbf{Severity:} Medium \textbf{Impact:} Inflexible governance
parameters

\textbf{Description:} The \texttt{vetoProposalPeriod} and
\texttt{vetoPeriod} are set at initialization and cannot be updated.
Changed governance requirements require contract redeployment.

\textbf{Recommendation:} Add admin functions to update parameters with
appropriate access control and timelock.

\begin{center}\rule{0.5\linewidth}{0.5pt}\end{center}

\subsubsection{M-05: Missing Vote Event in MultiDAOGovernor
Finalization}\label{m-05-missing-vote-event-in-multidaogovernor-finalization}

\textbf{Location:} \texttt{MultiDAOGovernor.sol:502-534}
(\texttt{finalizeProposal}) \textbf{Severity:} Medium \textbf{Impact:}
Missing audit trail

\textbf{Description:} No event is emitted when a proposal transitions to
\texttt{Succeeded} or \texttt{Defeated} state during finalization.

\textbf{Recommendation:} Add
\texttt{ProposalFinalized(bytes32\ indexed\ proposalId,\ ProposalState\ state)}
event.

\begin{center}\rule{0.5\linewidth}{0.5pt}\end{center}

\subsubsection{M-06: Unbounded Array Iteration in Charter Veto Voter
Management}\label{m-06-unbounded-array-iteration-in-charter-veto-voter-management}

\textbf{Location:} \texttt{Charter.sol:486-493}
(\texttt{removeVetoVoter}) \textbf{Severity:} Medium \textbf{Impact:}
Gas DoS

\textbf{Description:} Removing a veto voter iterates through the entire
\texttt{vetoVotersList} array. With many veto voters, this could exceed
block gas limits.

\textbf{Recommendation:} Use a mapping with linked list pattern or
maintain index mapping.

\begin{center}\rule{0.5\linewidth}{0.5pt}\end{center}

\subsubsection{M-07: Proposal State Race in
Council.sol}\label{m-07-proposal-state-race-in-councilsol}

\textbf{Location:} \texttt{Council.sol:180-217} (\texttt{proposalState})
\textbf{Severity:} Medium \textbf{Impact:} State inconsistency

\textbf{Description:} The proposal state is computed dynamically from
charter status and timestamps. Two calls to \texttt{proposalState()}
could return different values if block timestamp crosses a boundary
mid-transaction (in external calls).

\textbf{Recommendation:} Consider caching state or adding explicit state
transition functions.

\begin{center}\rule{0.5\linewidth}{0.5pt}\end{center}

\subsubsection{M-08: No Minimum Voting Period in
Charter.sol}\label{m-08-no-minimum-voting-period-in-chartersol}

\textbf{Location:} \texttt{Charter.sol:133-158} (\texttt{initialize})
\textbf{Severity:} Medium \textbf{Impact:} Flash governance attacks

\textbf{Description:} No minimum is enforced for
\texttt{votingPeriod\_}. A malicious deployer could set a 1-second
voting period, effectively allowing instant governance.

\textbf{Recommendation:} Add \texttt{MIN\_VOTING\_PERIOD} constant
(e.g., 1 day).

\begin{center}\rule{0.5\linewidth}{0.5pt}\end{center}

\subsubsection{M-09: Constitutional Amendment Bypass in
MultiDAOGovernor}\label{m-09-constitutional-amendment-bypass-in-multidaogovernor}

\textbf{Location:} \texttt{MultiDAOGovernor.sol:366-376}
\textbf{Severity:} Medium \textbf{Impact:} Governance capture

\textbf{Description:} The constitutional amendment check only verifies
array length and strategyDAO presence, not that the specified DAOs are
actually critical DAOs:

\begin{Shaded}
\begin{Highlighting}[]
\NormalTok{require(targetDAOs.length \textgreater{}= constitutionalApprovalCount, "Need more DAOs");}
\end{Highlighting}
\end{Shaded}

Attacker could use non-critical community DAOs to meet the count
requirement.

\textbf{Recommendation:} Verify critical DAO participation:

\begin{Shaded}
\begin{Highlighting}[]
\NormalTok{for (uint i = 0; i \textless{} constitutionalApprovalCount; i++) \{}
\NormalTok{    require(daos[targetDAOs[i]].isCritical, "Not a critical DAO");}
\NormalTok{\}}
\end{Highlighting}
\end{Shaded}

\begin{center}\rule{0.5\linewidth}{0.5pt}\end{center}

\subsection{Low Severity Findings}\label{low-severity-findings}

\subsubsection{L-01: Missing Zero Address Check in Council
Constructor}\label{l-01-missing-zero-address-check-in-council-constructor}

\textbf{Location:} \texttt{Council.sol:95-115} (\texttt{initialize})
\textbf{Impact:} Invalid initialization possible

\texttt{vault\_} and \texttt{target\_} are not validated for zero
address.

\begin{center}\rule{0.5\linewidth}{0.5pt}\end{center}

\subsubsection{L-02: Event Emission After State Change Pattern
Violation}\label{l-02-event-emission-after-state-change-pattern-violation}

\textbf{Location:} Multiple contracts \textbf{Impact:} Event ordering
issues

Events should be emitted after state changes. Several instances emit
before final state update.

\begin{center}\rule{0.5\linewidth}{0.5pt}\end{center}

\subsubsection{L-03: Floating Pragma in
MultiDAOGovernor}\label{l-03-floating-pragma-in-multidaogovernor}

\textbf{Location:} \texttt{MultiDAOGovernor.sol:4} \textbf{Impact:}
Compiler version inconsistency

Uses \texttt{\^{}0.8.24} while other contracts use \texttt{\^{}0.8.31}.

\begin{center}\rule{0.5\linewidth}{0.5pt}\end{center}

\subsubsection{L-04: Missing NatSpec
Documentation}\label{l-04-missing-natspec-documentation}

\textbf{Location:} Multiple internal functions \textbf{Impact:} Reduced
auditability

\begin{center}\rule{0.5\linewidth}{0.5pt}\end{center}

\subsubsection{L-05: Unused Error
Definitions}\label{l-05-unused-error-definitions}

\textbf{Location:} \texttt{Committee.sol} does not use custom errors
despite pattern elsewhere \textbf{Impact:} Inconsistent error handling

\begin{center}\rule{0.5\linewidth}{0.5pt}\end{center}

\subsubsection{L-06: Magic Numbers in Committee Block
Calculation}\label{l-06-magic-numbers-in-committee-block-calculation}

\textbf{Location:} \texttt{Committee.sol:131} \textbf{Impact:} Hardcoded
assumption

\begin{Shaded}
\begin{Highlighting}[]
\NormalTok{proposal.endBlock = block.number + votingPeriod / 12; // \textasciitilde{}12s blocks}
\end{Highlighting}
\end{Shaded}

Block time varies by chain.

\begin{center}\rule{0.5\linewidth}{0.5pt}\end{center}

\subsubsection{L-07: No Upgrade Gap in Upgradeable
Contracts}\label{l-07-no-upgrade-gap-in-upgradeable-contracts}

\textbf{Location:} \texttt{Charter.sol}, \texttt{Veto.sol},
\texttt{Council.sol} \textbf{Impact:} Future upgrade storage collision
risk

Contracts using EIP-7201 should still reserve gap for potential
non-namespaced additions.

\begin{center}\rule{0.5\linewidth}{0.5pt}\end{center}

\subsubsection{L-08: Redundant Interface Alias
Functions}\label{l-08-redundant-interface-alias-functions}

\textbf{Location:} \texttt{Charter.sol:537-603} \textbf{Impact:} Code
bloat, increased attack surface

Multiple alias functions (\texttt{charterAdmin()} -\textgreater{}
\texttt{admin()}, etc.) increase contract size unnecessarily.

\begin{center}\rule{0.5\linewidth}{0.5pt}\end{center}

\subsubsection{L-09: Missing Input Validation in MultiDAOGovernor
Constructor}\label{l-09-missing-input-validation-in-multidaogovernor-constructor}

\textbf{Location:} \texttt{MultiDAOGovernor.sol:204-236}
\textbf{Impact:} Invalid configuration possible

No validation that
\texttt{\_constitutionalApprovalCount\ \textless{}=\ \_criticalDAOs.length}.

\begin{center}\rule{0.5\linewidth}{0.5pt}\end{center}

\subsubsection{L-10: Secretariat setUp() Not
Protected}\label{l-10-secretariat-setup-not-protected}

\textbf{Location:} \texttt{Secretariat.sol:163} \textbf{Impact:}
Potential reinitialization

Abstract function requires implementers to ensure initialization
protection.

\begin{center}\rule{0.5\linewidth}{0.5pt}\end{center}

\subsubsection{L-11: VoteTrackerLRC20 Cannot Add Authorized
Callers}\label{l-11-votetrackerlrc20-cannot-add-authorized-callers}

\textbf{Location:} \texttt{VoteTrackerLRC20.sol:76-85} \textbf{Impact:}
Inflexible caller management

Once initialized, cannot add new authorized callers if governance
evolves.

\begin{center}\rule{0.5\linewidth}{0.5pt}\end{center}

\subsubsection{L-12: Committee Uses require() Instead of Custom
Errors}\label{l-12-committee-uses-require-instead-of-custom-errors}

\textbf{Location:} \texttt{Committee.sol:139-142,\ 163-167}
\textbf{Impact:} Higher gas costs, inconsistent error handling

\begin{center}\rule{0.5\linewidth}{0.5pt}\end{center}

\subsection{Informational}\label{informational}

\subsubsection{I-01: Governance.sol Contains Only
Re-exports}\label{i-01-governancesol-contains-only-re-exports}

The file only re-exports OpenZeppelin governance contracts. This is
acceptable but should be documented as a convenience module, not custom
implementation.

\subsubsection{I-02: Stake Token Division
Behavior}\label{i-02-stake-token-division-behavior}

The \texttt{Stake} contract inherits standard ERC20 division behavior.
Consider implications for vote weight calculations with tokens having
decimals != 18.

\subsubsection{I-03: ERC165 Implementation
Completeness}\label{i-03-erc165-implementation-completeness}

All contracts properly implement ERC165 for interface detection.

\subsubsection{I-04: Storage Slot Collision
Risk}\label{i-04-storage-slot-collision-risk}

EIP-7201 storage locations appear correctly computed, but should be
verified against keccak256 outputs.

\begin{center}\rule{0.5\linewidth}{0.5pt}\end{center}

\subsection{Recommendations Summary}\label{recommendations-summary}

\subsubsection{Immediate Actions
(Critical/High)}\label{immediate-actions-criticalhigh}

\begin{enumerate}
\def\labelenumi{\arabic{enumi}.}
\tightlist
\item
  \textbf{Add voting delay} between proposal creation and voting start
  (C-01)
\item
  \textbf{Implement vote weight tracking} across DAOs in
  MultiDAOGovernor (C-03)
\item
  \textbf{Add veto proposal cooldown} to prevent reset attacks (C-04)
\item
  \textbf{Verify attestations} in bridge tally reception (H-05)
\item
  \textbf{Make guard attachment immutable} or require parent approval
  (H-06)
\end{enumerate}

\subsubsection{Short-term Actions
(Medium)}\label{short-term-actions-medium}

\begin{enumerate}
\def\labelenumi{\arabic{enumi}.}
\setcounter{enumi}{5}
\tightlist
\item
  Implement quorum check in Committee.sol
\item
  Add minimum voting period constants
\item
  Emit finalization events in MultiDAOGovernor
\item
  Validate constitutional amendment DAO criticality
\end{enumerate}

\subsubsection{Long-term Actions
(Low/Informational)}\label{long-term-actions-lowinformational}

\begin{enumerate}
\def\labelenumi{\arabic{enumi}.}
\setcounter{enumi}{9}
\tightlist
\item
  Standardize error handling across all contracts
\item
  Add upgrade gaps to all upgradeable contracts
\item
  Remove redundant alias functions
\item
  Add comprehensive NatSpec documentation
\end{enumerate}

\begin{center}\rule{0.5\linewidth}{0.5pt}\end{center}

\subsection{Appendix: Attack Scenario
Details}\label{appendix-attack-scenario-details}

\subsubsection{Flash Loan Voting Attack (C-01
Detailed)}\label{flash-loan-voting-attack-c-01-detailed}

\begin{verbatim}
Block N:
1. Attacker calls Aave flash loan for 10M tokens
2. Attacker calls token.delegate(attacker)
3. Attacker calls council.submitProposal(...) via frontrunning
4. Charter.initializeProposal sets votingStartTimestamp = block.timestamp
5. Attacker calls charter.castVote with 10M token weight
6. Attacker repays flash loan

Result: Attacker controls vote outcome with zero capital
\end{verbatim}

\subsubsection{Cross-DAO Vote Amplification (C-03
Detailed)}\label{cross-dao-vote-amplification-c-03-detailed}

\begin{verbatim}
Setup: Attacker holds 1000 governance tokens
Proposal targets: [DAO_A, DAO_B, DAO_C, DAO_D, DAO_E]

1. castVote(proposalId, DAO_A, YES) -> 1000 weight to DAO_A
2. castVote(proposalId, DAO_B, YES) -> 1000 weight to DAO_B
3. castVote(proposalId, DAO_C, YES) -> 1000 weight to DAO_C
4. castVote(proposalId, DAO_D, YES) -> 1000 weight to DAO_D
5. castVote(proposalId, DAO_E, YES) -> 1000 weight to DAO_E

Total influence: 5000 weight from 1000 tokens (5x amplification)
\end{verbatim}

\begin{center}\rule{0.5\linewidth}{0.5pt}\end{center}

\subsection{Methodology}\label{methodology}

This audit employed:

\begin{itemize}
\tightlist
\item
  Manual code review
\item
  Control flow analysis
\item
  Data flow analysis
\item
  State machine modeling
\item
  Attack vector enumeration
\item
  Cross-contract interaction analysis
\end{itemize}

Automated tools used:

\begin{itemize}
\tightlist
\item
  Slither static analysis
\item
  Foundry fuzzing (where applicable)
\end{itemize}

\begin{center}\rule{0.5\linewidth}{0.5pt}\end{center}

\emph{This audit does not constitute legal or investment advice. Smart
contracts may contain vulnerabilities not identified in this report.}

\end{document}
